\section{Implementation and integration}
\label{sec:implementation}

\subsection{The order of operations}
Version 0.4.0 retains the preceding front-end structural certificates and
the existing simplification and invariant filters. Its inherited scanning
architecture follows the established local-complex approach~\cite{bnfast}. The standard scanner
remains the default. The new optional decision backend applies the following
operations after each crossing:
\begin{enumerate}
\item Form the exact crossing extension and use the inherited Gaussian
elimination to cancel unit maps.
\item For the Fitting backend, inspect bounded connected components, verify
candidate scalar splits, and transform the complete differential.
\item Recognize whole common-$\theta$ components and compute their interval
multiplicities from scalar product ranks.
\item In decision mode, shorten intervals to length at most two; in exact
mode, retain their lengths and absolute degree shifts.
\item Share identical interval or residual-component serializations and
update their nonnegative multiplicities. Decision multiplicities are capped
at three.
\end{enumerate}
The stage log records whether every surviving component was a singleton or
was actually recognized and normalized by the interval pass. A skipped
nontrivial component does not count as a certified checkpoint. The log also
records frontier size, unclassified object counts, and crossing distances
between checkpoints.

\subsection{Public interfaces}
The command line obtains a validated \code{Diagram} through the inherited
loader. The new complete decision routes are
\begin{lstlisting}
python -m fastunknot recognize examples/conway.json --backend barcode
python -m fastunknot recognize examples/conway.json --backend fitting
\end{lstlisting}
The usual inexpensive front-end checks may decide these examples before the
selected backend is needed. The reproduction scripts explicitly request raw
backend calculations when measuring scanner work.

\par\noindent\begin{minipage}{\linewidth}
For exact homology ranks, the separate routes are
\begin{lstlisting}
python -m fastunknot khovanov examples/conway.json --barcode
python -m fastunknot khovanov examples/conway.json --fitting
\end{lstlisting}
\end{minipage}\par
These retain complete interval lengths. The Python APIs are
\code{barcode\_khovanov\_rank}, \code{barcode\_khovanov\_decide},
\code{fitting\_khovanov\_rank}, and \code{fitting\_khovanov\_decide}, in the
corresponding modules. The Fitting exact-rank entry point rejects nontrivial
decision caps. Its optional \code{record\_witnesses=True} argument includes
split certificates in the result.

The optional ordering pass is selected, for example, by
\begin{lstlisting}
python -m fastunknot recognize examples/conway.json \
  --backend fitting --order-window 10 --order-passes 2
\end{lstlisting}
It optimizes overlapping windows using the frontier-mass objective. The
planner's time is included in the recognition budget. Standalone calls to
\code{improve\_scan\_order} expose both supported objectives and retain a
complete valid order if their optimization budget expires.

\subsection{Evidence and failure semantics}
An exact result contains \code{rank}, \code{reduced\_rank}, and
\code{by\_degree}. A default decision result contains \code{rank\_capped},
\code{rank\_cap=3}, and \code{length\_cap=2}. The latter does not report
the model's ordinary rank as the rank of the original knot.
The Python decision APIs also permit retaining full interval lengths or using
the universally safe length-three cutoff; these choices still return capped
rank evidence rather than an exact-rank schema.

Similarly, after length shortening, a weighted object count measures the
decision model, not an isomorphic expansion of the original complex.
The implementation changes the statistic name to
\code{weighted\_decision\_model\_objects} and sets a
\code{decision\_equivalence\_only} flag. This avoids incorrectly interpreting
an altered model as a lower bound on the original chain dimension.

The \code{max\_objects} ceiling is checked before constructing a crossing
extension, so a normalization that would shrink an oversized temporary model
is not performed if the ceiling would first be exceeded. Time budgets are
cooperative. The low-level APIs raise \code{ScanLimit} on exhaustion; the
recognition pipeline returns \UNKNOWN. Neither a large model nor a failed
decomposition search is treated as evidence of knottedness.

\subsection{Validation boundaries}
Checking $d^2=0$ is valuable, but it cannot establish observational equivalence
of two different complexes. The correctness of length shortening therefore
rests on Theorem~\ref{thm:geometric-cutoff}, while the tests independently
exercise its finite-matrix consequences. The optional \code{check\_d\_squared}
mode checks that subsequent algebraic operations still produce a complex.

The new normalization mixes quantum degrees because the inherited scanner
has already forgotten them. It preserves homological degree in exact mode;
its scalar Fitting blocks are explicitly restricted to matching and
homological degree. A future bigraded implementation must impose compatible
quantum shifts as an additional condition.

\subsection{Repository integration}
The package contains a complete runnable \code{fast/} tree, a reference copy
of the pinned repository files, and the preceding 0.3.0 reference used for
incremental comparison. One patch applies directly to the pinned
\code{Topology/UnknotRecognition/fast} tree; the other applies after the
structural-compression package. Both are generated from complete file bytes
and independently applied to fresh reference copies during packaging.

No remote repository modification is part of this deliverable. Integration
into a tree that also contains the independent three-braid continuation must
resolve shared front-end files and rerun the combined tests. The new algebra
modules and ordering module are otherwise self-contained, use only the
Python standard library, and retain the project's MIT-0 license.
